\section{Introduction}
\label{sec:intro}

LLMs are widely used as judges for preference annotation and post-training
\citep{bai2022constitutional,zheng2023judging,lee2024rlaif}. For a prompt $x$ and two responses
$y_A,y_B$, a judge queried by logit elicitation or repeated sampling returns a distribution
$\QJ(\cdot\mid X)$ over preference outcomes, $X=(x,y_A,y_B)$. This distribution is informative but
systematically different from the human preference we want to align to. Human labels are far more
expensive, so in practice we have judge outputs on many comparisons and human labels on only a
few. We ask how to combine the two into a human-anchored preference target, and then how to train a
policy against that target.

\paragraph{Starting point.}
\citet{li2026personalizing} study few-shot personalization of a black-box predictor in
nonparametric regression. Their procedure has three steps:
\begin{enumerate}
  \item apply a local smoothing to the pretrained predictor;
  \item estimate the discrepancy between the target and the smoothed predictor with a kernel
        smoother on a few target labels;
  \item choose the amount of smoothing on validation data.
\end{enumerate}
The candidate set in step~3 includes the target-only estimator, so an uninformative predictor is
automatically ignored. They also give a variance-driven rule for where to collect labels, and prove
a minimax rate that improves on target-only estimation when the discrepancy is simpler than the
target. Their setting is scalar regression, which includes binary classification. It does not
cover graded black-box outputs or preference post-training.

\paragraph{This proposal.}
We extend this idea in three ways.
\begin{enumerate}
  \item \textbf{Distributional few-shot personalization (\dfsp).} The judge's win probability is the
        offset. The discrepancy is localized on richer summaries of a graded judge distribution
        (tie mass, order inconsistency). This lets the correction use information that a scalar
        probability discards.
  \item \textbf{Label acquisition.} A design rule for which comparisons to send to humans, derived
        from the risk of the estimator. It favors regions where humans disagree and where the
        judge-to-human discrepancy changes quickly.
  \item \textbf{Post-training (\fsppo).} The corrected probabilities are soft targets for DPO-style
        training. We show how the estimation error bounds the excess preference risk of the
        trained policy.
\end{enumerate}

\paragraph{Related work.}
Soft AI preference labels are already used for reward modeling and preference optimization
\citep{lee2024rlaif,furuta2024gapo}, as are strength-aware and distributional preference models
\citep{imai2026mixdpo,li2024dprm}. Other work combines few human labels with AI feedback through
targeted correction \citep{xu2025rlthf}, a learned linear map from judge outputs to human judgments
\citep{vandenburg2025aligning}, or a jointly weighted Bradley--Terry likelihood whose weight can fall
back to human-only estimation \citep{cai2026dial}. Active selection of pairs for human labeling has
also been studied \citep{shen2025active}. Our contribution is none of these components on its own.
It is a nonparametric, adaptive correction of a \emph{distribution-valued} judge, with guarantees
that carry through to post-training.
